\chapter{Introducción}
\label{cap:intro}

En este capítulo se establece el marco general del trabajo, definiendo la relevancia de la computación cuántica actual y el problema crítico del ruido.

\section{Contexto: La era NISQ}
Descripción de los dispositivos \textit{Noisy Intermediate-Scale Quantum} y por qué la computación cuántica actual depende de la mitigación de errores.

\section{Motivación y Justificación}
Por qué el ruido es el principal obstáculo para la escalabilidad. El valor de aplicar Inteligencia Artificial frente a métodos puramente estadísticos.

\section{Objetivos del Trabajo}
\subsection{Objetivo Principal}
Desarrollar un pipeline híbrido de IA para la mitigación de errores de puerta y lectura.
\subsection{Objetivos Específicos}
Implementación de Transformers (GEM), GNNs (REM) y validación en hardware real de IBM.

\section{Metodología y Estructura}
Breve resumen de las fases del TFM (Ingeniería de datos, Modelado y Validación).